\section{三个二次的关系}

一元二次函数、一元二次方程、一元二次不等式是同一件事的三种说法。
看一张图就够了：

\begin{figure}[htbp]
  \centering
  \includegraphics[width=0.62\textwidth]{parabola.pdf}
  \caption{$y=x^2-2x-3$：交点、上方、下方对应三种问题}
\end{figure}

\begin{itemize}
  \item \key{方程}：令 $y=0$，求与 $x$ 轴交点的横坐标，得 $x=-1$ 或 $x=3$；
  \item \key{大于零}：$y>0$ 的部分在 $x$ 轴上方，解集为 $(-\infty,-1)\cup(3,+\infty)$；
  \item \key{小于零}：$y<0$ 的部分在 $x$ 轴下方，解集为 $(-1,3)$。
\end{itemize}

\key{一句话记忆}：方程找交点，不等式找上下，根是边界。

\section{数轴穿根法}

数轴穿根法用来解\key{高次不等式}和\key{分式不等式}，比逐段讨论更快。
它回答的问题是：\key{每个区间上，式子的符号是正还是负}。

\subsection{先处理最高次项的系数}

穿根前必须先看清最高次项系数的符号，它决定从哪一端起笔：

\begin{center}
\footnotesize
\setlength{\tabcolsep}{3pt}
\renewcommand{\arraystretch}{1.4}
\begin{tabular}{p{4.2cm}p{4.6cm}p{4.2cm}}
\toprule
最高次项系数 & 起笔位置 & 做法 \\
\midrule
为正 & 数轴右上方 & 直接穿根 \\
为负 & 数轴右下方 & 两边同乘 $-1$，不等号方向改变，化为正 \\
\bottomrule
\end{tabular}
\end{center}

\key{一句话}：\textbf{一律先化为正}，再从右上方起笔。
这样规则只有一条，不必记两套。

\subsection{操作步骤}

\begin{enumerate}
  \item 把不等式化成「左边是因式乘积，右边是 $0$」，并使最高次项系数为正；
  \item 求出所有根，按从小到大标在数轴上；
  \item 从数轴右上方起笔，依次穿过每个根；
  \item \key{奇穿偶不穿}：根的重数是奇数（奇重根）就穿过数轴，是偶数（偶重根）就弹回；
  \item 按不等号方向取数轴上方（$>0$）或下方（$<0$）的区间；
  \item 检查端点：带等号时包含根。
\end{enumerate}

\begin{example}
解不等式 $(x+1)(x-2)(x-4)<0$。
\end{example}
\begin{solution}
最高次项系数为正。根为 $-1$、$2$、$4$（都是一重根，全部穿过）。
从右上方起笔，依次穿过 $4$、$2$、$-1$，得到符号分布：
$(-\infty,-1)$ 为负，$(-1,2)$ 为正，$(2,4)$ 为负，$(4,+\infty)$ 为正。
取小于 $0$ 的部分，解集为 $(-\infty,-1)\cup(2,4)$。
\end{solution}

\begin{figure}[htbp]
  \centering
  \includegraphics[width=0.70\textwidth]{ineq-on-axis.pdf}
  \caption{穿根折线与解集（空心点表示端点不取）}
\end{figure}

\key{一句话记忆}：右上方起笔，奇穿偶不穿，大于取上、小于取下。

\section{分式不等式的解法}

分式不等式用同一条思路：\key{先判断正负，再画图求解}。

分式 $\dfrac{f(x)}{g(x)}$ 的符号由分子、分母共同决定：
\[
  \frac{f(x)}{g(x)}>0 \iff f(x)g(x)>0 \quad (g(x)\neq0).
\]
于是分式不等式可以\key{转化为整式不等式}，再用数轴穿根法。

\key{特别注意}：分母的零点永远取不到（它使分母为 $0$），必须画\key{空心点}。

\begin{example}
解不等式 $\dfrac{x+1}{x-2}\geq0$。
\end{example}
\begin{solution}
等价于 $(x+1)(x-2)\geq0$ 且 $x\neq2$。
根为 $-1$、$2$，最高次项系数为正。
从右上方起笔依次穿根，得符号：$(-\infty,-1]$ 为正，$(-1,2)$ 为负，$(2,+\infty)$ 为正。
取大于等于 $0$ 的部分，并去掉分母零点 $2$，
解集为 $(-\infty,-1]\cup(2,+\infty)$。
\end{solution}

\begin{figure}[htbp]
  \centering
  \includegraphics[width=0.70\textwidth]{rational-ineq.pdf}
  \caption{$\dfrac{x+1}{x-2}\geq0$：$-1$ 实心（可取），$2$ 空心（分母零点）}
\end{figure}

\section{解集在数轴上的表示}

这一节与前面不同：穿根法是\key{求}解集的方法，
本节是把已经求出的解集\key{画}出来。

\begin{figure}[htbp]
  \centering
  \includegraphics[width=0.56\textwidth]{interval-notation.pdf}
  \caption{区间在数轴上的三种端点情形}
\end{figure}

\key{约定}：空心圈表示不含该端点，实心点表示含有该端点。
例如 $x^2-2x-3<0$ 的解集 $(-1,3)$ 画成两个空心圈；
$x^2-2x-3\leq0$ 的解集 $[-1,3]$ 画成两个实心点。

\section{用数轴表示集合与条件关系}

判断条件关系时，把 $p$、$q$ 的解集画在数轴上，包含关系一目了然。
建议按「先认端点 $\to$ 再看包含」的顺序看下图：

\begin{figure}[htbp]
  \centering
  \includegraphics[width=0.70\textwidth]{set-conditions.pdf}
  \caption{四种条件关系的数轴表示（分步）}
\end{figure}

\key{一句话记忆}：解集画在数轴上，小范围推大范围。

\section{温故题卡}

每课开头用 5 分钟做 2--3 道题，复现上一课的结论，不重新讲解。

\begin{table}[htbp]
\centering
\footnotesize
\setlength{\tabcolsep}{3pt}
\renewcommand{\arraystretch}{1.4}
\begin{tabular}{p{1.6cm}p{9.8cm}p{2.2cm}}
\toprule
\key{课次} & \key{温故题} & \key{复现的结论} \\
\midrule
第 2 课 & 把 $y=2x^2-4x+1$ 配方成顶点式 & 配方法 \\
第 3 课 & 求 $y=x^2-2x-3$ 的顶点与最值 & 顶点与最值 \\
第 4 课 & 解不等式 $x^2-2x-3<0$ & 解不等式流程 \\
\bottomrule
\end{tabular}
\caption{温故题卡}
\end{table}
